\begin{abstract}
\textbf{Objective:}
Regularisation strength selected for covariance fit may differ from
the strength selected for discrimination of unseen EEG participants.
We investigate the consequence for zero-target cross-subject
decoding and the validity of spectral-representation comparisons,
while auditing the precise source-selection objectives.
\textbf{Methods:}
We compare participant-grouped source macro-AUC selection with a
class-centred held-source covariance Gaussian log score
(the original runner's historically named covariance-NLL criterion)
and analytic Ledoit--Wolf in a frozen covariance/tangent-LDA
pipeline. Independent EEG cohorts, paired participant inference,
prospectively frozen representation controls and original frozen
trial-score audits are reported separately. We distinguish the
original covariance log score from a genuinely predictive
class-conditional Gaussian log score, which is proposed but
has not yet been tested against EEG outcomes.
\textbf{Results:}
On EEGMMIDB S021--S040 (\(N=20\)), source-macro selection
increased mean AUC by 0.0561 over the class-centred covariance
score selector (Holm \(p=0.00597\)); five participants were
harmed by more than 0.02 AUC. Two independently frozen BNCI
cohorts had positive average directions but individually
non-significant adjusted comparisons. Separately frozen
Shin2017A (\(N=29\)) reported +0.0222 AUC (Holm \(p=0.0223\));
the original score archive passed independent checksum
verification. An additional prospective filter-bank experiment
contradicted the originally hypothesised universal advantage.
A new estimand audit proves that class-specific translation
of held-source validation features leaves the frozen
covariance score unchanged, unlike proper predictive likelihood.
\textbf{Conclusions:}
Objective mismatch has experimentally measurable consequences
in the studied source-only EEG setting, but results against
a validation-class-centred covariance score must not be
misrepresented as results against true out-of-sample
predictive likelihood. No universal shrinkage strength,
target-safe switching rule or population-level spectral
advantage is established.
\end{abstract}
